\section{Introduction}
\label{sec:intro}

Modern deep learning is characterized by an explosion of task-specific models. However, maintaining separate specialized networks is impractical in resource-constrained edge environments---such as microcontrollers, wearables, and IoT devices---where memory, storage, and compute budgets are severely restricted. Model merging integrates multiple task-specific experts (e.g., LoRA adapters \cite{lora}) into a single shared backbone \cite{merging_survey}, dramatically reducing storage and parameter-management overhead.

While standard model merging is evaluated almost exclusively under expensive full-precision (FP16/FP32) regimes, real-world edge devices require post-training quantization (PTQ) to low-bit integer formats (INT4/INT8) to avoid prohibitive latency and battery drain \cite{gholami2021survey}. When deploying multi-task merged models to the edge, two severe bottlenecks arise: (1) \emph{Linear Scaling of Ensembling:} Parallel execution of $K$ experts in activation-space scales compute and SRAM linearly ($O(K)$), which is prohibitively expensive. (2) \emph{Static vs. Dynamic Modularity:} Weight-space merging (Post-Merge Quantization or PMQ) runs in constant $O(1)$ compute but collapses catastrophically under realistic non-linear activations (e.g., GELU) due to representation misalignment. Furthermore, static merging lacks modularity: registering new tasks requires complete offline re-merging and re-quantization.

\begin{table}[t]
\caption{Architectural and physical serve-time trade-offs on resource-constrained edge devices.}
\label{tab:tradeoffs}
\vskip 0.05in
\begin{center}
\begin{small}
\begin{sc}
\setlength{\tabcolsep}{3pt}
\begin{tabular}{lcccc}
\toprule
Paradigm & Compute & SRAM & Modular & Accuracy \\
\midrule
Ensembling & $O(K)$ & High & High & High \\
PMQ & $O(1)$ & Zero & Low & Collapse \\
\textbf{SA-QAB (Ours)} & \textbf{$O(1)$} & \textbf{Low} & \textbf{High} & \textbf{Robust} \\
\bottomrule
\end{tabular}
\end{sc}
\end{small}
\end{center}
\vskip -0.1in
\end{table}

To resolve these pressing practical challenges, we propose \textbf{Scale-Aligned Quantized Activation Blending (SA-QAB)}, a training-free, forward-pass-only framework designed for robust, multi-task edge execution. SA-QAB keeps the heavy base model and specialized task adapters separated during quantization and executes them in their native integer formats. Crucially, by utilizing a low temperature during dynamic routing, SA-QAB achieves near-sparse routing, allowing the hardware to execute \emph{only the single active expert pathway} (constant $O(1)$ adapter compute) instead of all $K$ pathways in parallel, while maintaining a dynamic decoupled serving structure that allows on-the-fly task loading and registration.

Our primary contributions are:
(1) \textbf{Decoupled Heterogeneous Quantization (DHQ):} We propose a hardware-aligned strategy compressing heavy base weights to INT4 to minimize memory latency, while keeping lightweight LoRA expert weights in INT8 to preserve task-specific representation boundaries.
(2) \textbf{Quantization Scale Recovery (QSR):} We introduce a training-free calibration technique pre-computing scale recovery factors on a 64-sample set to correct for low-bit scale contraction on-the-fly and restore representation fidelity.
(3) \textbf{Quantized Zero-Shot Centroid Alignment (Q-ZCA):} We implement an early-stage dynamic routing layer operating in integer space, comparing quantized activations with task centroids using INT8 cosine similarity to route inputs sample-by-sample.
(4) \textbf{Evaluation and Systems Trade-off Analysis:} We evaluate SA-QAB on a 192D simulation suite and real pixels on ViT-Tiny across MNIST, Fashion-MNIST, CIFAR-10, and SVHN under non-linear (GELU) activations. While static merging collapses to \textbf{18.60\%} accuracy, SA-QAB achieves \textbf{77.50\%} joint accuracy (a \textbf{+58.90\% absolute improvement}). We also present GMM and SRAM edge trade-offs.
